\chapter{Quantum Computation Basics}
% Solutions to the exercises of \cref{ch:lec01}.

\begin{solution}[Exercise~\ref{ex:lec01-1}]
$\norm{\psi}^{2}=\abs{2\ii}^{2}+\abs{-1}^{2}=4+1=5$, so $\norm{\psi}=\sqrt5$. The
normalized state is $\ket{\hat\psi}=\tfrac{1}{\sqrt5}(2\ii\ket{0}-\ket{1})$. The
probabilities are $\abs{2\ii/\sqrt5}^{2}=4/5$ for outcome $0$ and
$\abs{-1/\sqrt5}^{2}=1/5$ for outcome $1$; they sum to $1$ as required.
\end{solution}

\begin{solution}[Exercise~\ref{ex:lec01-2}]
$\braket{+|+}=\tfrac12(1+1)=1$ and $\braket{-|-}=\tfrac12(1+1)=1$, so both are unit
vectors. $\braket{+|-}=\tfrac12(\braket{0|0}-\braket{0|1}+\braket{1|0}-\braket{1|1})=\tfrac12(1-0+0-1)=0$,
so they are orthogonal; two orthonormal vectors in a $2$-dimensional space form a
basis. Adding and subtracting, $\ket{0}=\tfrac{1}{\sqrt2}(\ket{+}+\ket{-})$ and
$\ket{1}=\tfrac{1}{\sqrt2}(\ket{+}-\ket{-})$.
\end{solution}

\begin{solution}[Exercise~\ref{ex:lec01-3}]
For unit vectors, $\norm{u}=\norm{v}=1$, so Cauchy--Schwarz gives
$\abs{\braket{u|v}}\le1$ with equality iff $\ket{u},\ket{v}$ are linearly dependent,
i.e. $\ket{v}=\lambda\ket{u}$ for some $\lambda\in\C$. Taking norms,
$1=\norm{v}=\abs{\lambda}\norm{u}=\abs{\lambda}$, so $\lambda=\ee^{\ii\gamma}$ for some
real $\gamma$. Conversely, if $\ket{v}=\ee^{\ii\gamma}\ket{u}$ then
$\braket{u|v}=\ee^{\ii\gamma}$ and $\abs{\braket{u|v}}=1$. By the invariance of
physical states under a global phase, the two are the same physical state.
\end{solution}

\begin{solution}[Exercise~\ref{ex:lec01-4}]
$P^{2}=\ket{u}\braket{u|u}\bra{u}=\proj{u}=P$ since $\braket{u|u}=1$. Taking the
adjoint, $P\dg=(\proj{u})\dg=\proj{u}=P$. For the trace, using the identity
$\Tr(\ket{u}\!\bra{v})=\braket{v|u}$ we get $\Tr(P)=\braket{u|u}=1$ (a rank-one
projector has trace equal to its rank).
\end{solution}

\begin{solution}[Exercise~\ref{ex:lec01-5}]
Compare with the Bloch form
$\ket{\psi}=\cos\tfrac{\theta}{2}\ket{0}+\ee^{\ii\varphi}\sin\tfrac{\theta}{2}\ket{1}$.
The $\ket{0}$ amplitude is real and nonnegative, $\cos\tfrac{\theta}{2}=\tfrac12$,
giving $\theta/2=\pi/3$, so $\theta=2\pi/3$. Then
$\sin\tfrac{\theta}{2}=\sin\tfrac{\pi}{3}=\tfrac{\sqrt3}{2}$, matching the modulus of
the $\ket{1}$ amplitude, and its phase is $\varphi=\pi/3$. Hence
$(\theta,\varphi)=(2\pi/3,\ \pi/3)$.
\end{solution}

\begin{solution}[Exercise~\ref{ex:lec01-6}]
Factor out the first qubit:
$\ket{\psi}=\tfrac{1}{\sqrt2}\ket{0}\otimes(\ket{0}+\ket{1})
=\ket{0}\otimes\tfrac{1}{\sqrt2}(\ket{0}+\ket{1})=\ket{0}\otimes\ket{+}$.
This is an explicit product of single-qubit states, so $\ket{\psi}$ is separable (a
product state), not entangled.
\end{solution}

\begin{solution}[Exercise~\ref{ex:lec01-7}]
$H$ has real symmetric entries, so $H\dg=H^{\top}=H$. Then
$H^{2}=\tfrac12\left(\begin{smallmatrix}1&1\\1&-1\end{smallmatrix}\right)
\left(\begin{smallmatrix}1&1\\1&-1\end{smallmatrix}\right)
=\tfrac12\left(\begin{smallmatrix}2&0\\0&2\end{smallmatrix}\right)=I$, which shows
$H\dg H=H^{2}=I$, so $H$ is unitary and self-inverse. For $HYH$, use $Y=\ii XZ$
(\cref{prop:lec01-yxz}) and $H^{2}=I$:
\[
  HYH=\ii\,H X Z H = \ii(HXH)(HZH)=\ii\,Z\,X
     = \ii\begin{pmatrix}1&0\\0&-1\end{pmatrix}
       \begin{pmatrix}0&1\\1&0\end{pmatrix}
     = \ii\begin{pmatrix}0&1\\-1&0\end{pmatrix}
     = \begin{pmatrix}0&\ii\\-\ii&0\end{pmatrix}=-Y.
\]
So $HYH=-Y$.
\end{solution}

\begin{solution}[Exercise~\ref{ex:lec01-8}]
Write $\ket{+}\ket{+}=\tfrac12(\ket{00}+\ket{01}+\ket{10}+\ket{11})$. Apply $\CNOT$
(control $0$) term by term: $\ket{00}\to\ket{00}$, $\ket{01}\to\ket{01}$,
$\ket{10}\to\ket{11}$, $\ket{11}\to\ket{10}$. The sum is invariant, so
$\CNOT\ket{+}\ket{+}=\tfrac12(\ket{00}+\ket{01}+\ket{11}+\ket{10})=\ket{+}\ket{+}$.
The result is a product state, hence \emph{not} entangled: $\ket{+}\ket{+}$ is a
$+1$ eigenstate of this $\CNOT$.

For the cloning of orthogonal states, take $\ket{s}=\ket{0}$ and $U=\CNOT$ with
control $0$, target $1$. Then $\CNOT\ket{0}\ket{0}=\ket{0}\ket{0}$ and
$\CNOT\ket{1}\ket{0}=\ket{1}\ket{1}$, so both orthogonal basis states are copied.
This does not contradict \cref{thm:lec01-nocloning}: the theorem only forbids a
\emph{single} gate that copies \emph{every} state. This $\CNOT$ copies
$\ket{0},\ket{1}$ but fails to copy a superposition $a\ket{0}+b\ket{1}$, producing the
entangled $a\ket{00}+b\ket{11}$ instead. Cloning a fixed orthonormal set is exactly
the classical copy the theorem permits.
\end{solution}

\begin{solution}[Exercise~\ref{ex:lec01-9}]
(a) With $\rho=\proj{\psi}=\tfrac12(\proj{00}+\ket{00}\!\bra{11}+\ket{11}\!\bra{00}+\proj{11})$,
tracing out the second qubit and keeping only terms whose second-qubit bra and ket
labels agree gives
\[
  \rho_A=\Tr_B[\rho]=\tfrac12\big(\proj{0}\braket{0|0}+\proj{1}\braket{1|1}\big)
        =\tfrac12(\proj{0}+\proj{1})=\tfrac12\id.
\]
Its purity is $\Tr[\rho_A^{2}]=\Tr[\tfrac14\id]=\tfrac14\cdot2=\tfrac12<1$, so the
reduced state is maximally mixed.

(b) The projector for outcome $0$ on the first qubit is $\proj{0}\otimes\id$. The
outcome probability equals the squared norm of the unnormalized post-measurement
vector $(\proj{0}\otimes\id)\ket{\psi}=\tfrac{1}{\sqrt2}\ket{00}$, namely
$\bigl\|\tfrac{1}{\sqrt2}\ket{00}\bigr\|^{2}=\tfrac12$. Equivalently
$p(0)=\Tr[(\proj{0}\otimes\id)\rho]=\tfrac12$, consistent with $\rho_A=\tfrac12\id$.
After normalization the state is $\ket{0}\otimes\ket{0}$, so the second qubit is left
in $\ket{0}$---perfectly correlated with the first outcome.

For comparison, the product state $\ket{+}\ket{+}$ has reduced first-qubit state
$\rho_A'=\proj{+}$, a \emph{pure} state with purity $\Tr[\rho_A'^{2}]=1$: no
information is hidden in correlations because the state is separable. The contrast
$\tfrac12<1$ versus $1$ quantifies the entanglement of the Bell pair.
\end{solution}

\begin{solution}[Exercise~\ref{ex:lec01-10}]
(a) From the grouped expression \eqref{eq:lec01-tele-grouped}, the $(m_1,m_2)=(1,0)$
branch carries qubit~$3$ in the state $a\ket{0}-b\ket{1}$. Since $Z\ket{0}=\ket{0}$ and
$Z\ket{1}=-\ket{1}$, indeed $Z\ket{\psi}=Z(a\ket{0}+b\ket{1})=a\ket{0}-b\ket{1}$, so
Bob holds $Z\ket{\psi}$.
(b) The correction is $Z^{m_1}X^{m_2}=Z^{1}X^{0}=Z$. Applying it,
$Z\,(Z\ket{\psi})=Z^2\ket{\psi}=\ket{\psi}$ because $Z^2=I$.
(c) Averaging over the four outcomes, Bob's ensemble is
$\{\tfrac14,\ket{\psi}\},\{\tfrac14,X\ket{\psi}\},\{\tfrac14,Z\ket{\psi}\},\{\tfrac14,XZ\ket{\psi}\}$,
with density operator $\rho_3=\tfrac14\bigl(\proj{\psi}+X\proj{\psi}X\dg+Z\proj{\psi}Z\dg+XZ\proj{\psi}(XZ)\dg\bigr)$.
Writing $\ket{\psi}\!\bra{\psi}=\tfrac12(\id+\vec r\cdot\vec\sigma)$ in the Pauli basis,
the four Pauli conjugations flip the signs of the $\vec r\cdot\vec\sigma$ term so that
the four Bloch vectors sum to zero, leaving $\rho_3=\tfrac12\id$. Because this is
independent of $a,b$ and of any operation Alice performs, Bob's local statistics cannot
convey information until the classical bits $(m_1,m_2)$ arrive over a light-speed-limited
channel---hence no superluminal signaling.
\end{solution}

\begin{solution}[Exercise~\ref{ex:lec01-11}]
Encoding $(x,y)=(1,1)$: apply $ZX$ to qubit~$1$ of
$\ket{\Phi^{+}}=\tfrac{1}{\sqrt2}(\ket{00}+\ket{11})$. With the stated convention
$ZX\ket{0}=\ket{1}$ and $ZX\ket{1}=-\ket{0}$,
\[
  (ZX\otimes I)\ket{\Phi^{+}}=\tfrac{1}{\sqrt2}\bigl(\ket{10}-\ket{01}\bigr).
\]
Now Bob applies $\CNOT_{1\to2}$ (flip qubit~$2$ when qubit~$1$ is $\ket{1}$):
$\ket{10}\mapsto\ket{11}$, $\ket{01}\mapsto\ket{01}$, giving
$\tfrac{1}{\sqrt2}(\ket{11}-\ket{01})=\tfrac{1}{\sqrt2}(\ket{1}-\ket{0})_1\otimes\ket{1}_2$.
Applying $H$ to qubit~$1$ and using $H\tfrac{1}{\sqrt2}(\ket{1}-\ket{0})=H\bigl(-\ket{-}\bigr)=-\ket{1}$,
the state becomes $-\ket{11}$. A computational-basis measurement therefore returns
$11$ with probability $1$ (the global sign $-1$ is unobservable), recovering
$(x,y)=(1,1)$.
\end{solution}

\begin{solution}[Exercise~\ref{ex:lec01-12}]
(a) The four differences are $\alpha_1-\beta_1=-\tfrac{\pi}{4}$,
$\alpha_1-\beta_2=-\tfrac{3\pi}{4}$, $\alpha_2-\beta_1=\tfrac{\pi}{4}$,
$\alpha_2-\beta_2=-\tfrac{\pi}{4}$. Hence
$\E[S]=\cos(-\tfrac{\pi}{4})-\cos(-\tfrac{3\pi}{4})+\cos(\tfrac{\pi}{4})+\cos(-\tfrac{\pi}{4})
=\tfrac{1}{\sqrt2}+\tfrac{1}{\sqrt2}+\tfrac{1}{\sqrt2}+\tfrac{1}{\sqrt2}=2\sqrt2$.
(b) Fix $\beta_1,\beta_2$ and group the $\alpha_1$- and $\alpha_2$-dependent terms:
\[
  f=\bigl[\cos(\alpha_1-\beta_1)-\cos(\alpha_1-\beta_2)\bigr]
    +\bigl[\cos(\alpha_2-\beta_1)+\cos(\alpha_2-\beta_2)\bigr].
\]
Each bracket has the form $R\cos(\alpha-\phi)$ for suitable $R\ge0,\phi$: writing
$c\cos\alpha+s\sin\alpha=R\cos(\alpha-\phi)$ with $R=\sqrt{c^2+s^2}$, the first bracket
has amplitude $R_1=\sqrt{2-2\cos(\beta_1-\beta_2)}$ and the second
$R_2=\sqrt{2+2\cos(\beta_1-\beta_2)}$. Optimizing over $\alpha_1,\alpha_2$ independently
sets each cosine to $1$, giving $\max_{\alpha_1,\alpha_2}f=R_1+R_2$. Then
\[
  (R_1+R_2)^2=R_1^2+R_2^2+2R_1R_2=4+2\sqrt{(2-2c)(2+2c)}=4+4\sqrt{1-c^2}\le 8,
\]
where $c=\cos(\beta_1-\beta_2)\in[-1,1]$, with equality iff $c=0$, i.e.
$\beta_1-\beta_2=\pm\tfrac{\pi}{2}$. Thus $\max f=\sqrt8=2\sqrt2$, so
$\abs{\E[S]}\le 2\sqrt2$ for this state and the chosen settings attain it.
\end{solution}
